% Propietario conservado: /Users/ruben/Documents/ChatGPT/jueces y controles/REVISION_ARGUMENTAL_APERTURAS_CIERRES_20260915/II/manuscript_es/sections/10_gravedad.tex
\section{Acción, gravedad e información térmica}
\label{ii:sec:gravity}
\subsection{Período dodecafásico y realización circular de la acción}
El calendario ya construido determina $T_{108}=108t_0$.
En su realización circular,
\[
\omega_{108}=\frac{2\pi}{108t_0},\qquad
R_{108}=\frac{c_{\rm int}}{\omega_{108}}
=\frac{54}{\pi}\ell_0,\qquad \ell_0=c_{\rm int}t_0.
\]
La unidad $\ell_0$ y las cartas de tiempo y velocidad se mantienen
explícitas; la carta SI no selecciona el ciclo.
Para cada sección orientada de acción $a\in\{\mathrm{pre},\mathrm{ret}\}$,
\[
E_{108,a}=\hbar_a\omega_{108},\qquad
m_{108,a}=\frac{E_{108,a}}{c_{\rm int}^2}.
\]
La longitud reducida y la completa son respectivamente
\[
\frac{\hbar_a}{m_{108,a}c_{\rm int}}=R_{108},
\qquad
\frac{h_a}{m_{108,a}c_{\rm int}}=2\pi R_{108}.
\]
La misma acción y la misma duración intervienen en ambas lecturas. Su coincidencia radial se determina en la proposición siguiente.

\begin{proposition}[Selector de coincidencia de radios]
En la familia dimensional
\[
G_a(\vartheta)=\vartheta\frac{c_{\rm int}^5t_0^2}{\hbar_a},
\qquad \vartheta>0,
\]
la condición de la realización circular
$G_a(\vartheta)m_{108,a}/c_{\rm int}^2=R_{108}$
tiene la única solución
\[
\vartheta^\circ_{108}=(54/\pi)^2.
\]
\end{proposition}
\begin{proof}
El primer radio es
$r_g(\vartheta)=\vartheta(\pi/54)\ell_0$.
Compararlo con $(54/\pi)\ell_0$, con $\ell_0>0$, reduce la
igualdad a una ecuación lineal estrictamente creciente.
\end{proof}
La condición de coincidencia es parte explícita de esta realización;
el teorema resuelve su coeficiente sin usar un valor metrológico de $G$.
El radio aquí empleado es $Gm/c^2$, por lo que no se traslada
inadvertidamente el factor dos de la convención $2Gm/c^2$.

\begin{proposition}[Acción y gravedad recíprocas, área conservada]
En las dos cartas anteriores,
\[
\frac{G_{\rm pre}}{G_{\rm ret}}=R_{\rm act}^{-1},\qquad
\hbar_aG_a=\vartheta^\circ_{108}c_{\rm int}^5t_0^2,
\qquad
\ell_P^2=\frac{\hbar_aG_a}{c_{\rm int}^3}
=\vartheta^\circ_{108}\ell_0^2.
\]
\end{proposition}
\begin{proof}
La definición de $G_a$ contiene $\hbar_a^{-1}$ y el resto de factores
es común a las dos cartas. Tomar el cociente y el producto da
las tres igualdades.
\end{proof}
Esta reciprocidad explica por qué la información de orientación puede
variar sin que cambie el área. Se dispone así de una relación precisa
entre la elipse de acción, la escala de Planck y el operador areal
de la sección~\ref{ii:sec:area}.

\subsection{Normalización térmica y energía por ciclo}
Fijamos la misma sección orientada \(a\) para la energía y la
temperatura. La relación térmica del corpus publica
\begin{equation}
k_B\Theta_{{\rm clk},a}\log3
=\frac{h_a}{108t_0}
=\frac{\hbar_a}{t_P}=E_{P,a},
\qquad t_P=\frac{54}{\pi}t_0.
\label{ii:eq:thermalclock}
\end{equation}
El factor $\log3$ mide la información de la terna equiprobable;
la energía del reloj incorpora la acción ya generada.
La identidad conserva el producto $k_B\Theta_{\rm clk}$ como sección
de energía y hace explícita la carta térmica utilizada, desarrollada en la sección~\ref{ii:sec:ii-kb-aut-desarrollo}.

La realización radiativa proporciona, en esta escala,
\[
n_\gamma\ell_P^3=\frac{2\zeta(3)}{\pi^2\log^3 3},\qquad
\frac{u_\gamma\ell_P^3}{E_P}
=\frac{\pi^2}{15\log^4 3},\qquad
\frac{s_\gamma\ell_P^3}{k_B}
=\frac{4\pi^2}{45\log^3 3}.
\]
Aquí se encuentran dos invariantes distintos: $\log3$ conserva la
información TRIT y $\zeta(3)$ el momento de ocupación bosónica.
Su presencia en la misma expresión relaciona las constantes térmicas
con los momentos espectrales de la ocupación bosónica, sin identificar
estas magnitudes entre sí.

La sección~\ref{ii:sec:ii-kb-aut-desarrollo} desarrolla el transductor
de Boltzmann, su acción sobre la información trítica y el cambio de
coordenadas térmicas. La ecuación~\eqref{ii:eq:thermalclock} fija el
producto energía--temperatura que utiliza esa construcción. La
selección de una coordenada individual se distingue allí de la
identidad para el producto. Las fórmulas radiativas anteriores
conservan, en este mismo desarrollo, el contenido de ocupación y
entropía de la realización bosónica.

\subsection{Barbero--Immirzi en la realización Cartan--Holst}
Una vez construido $\gamma_{\HMT}$, la fuente lo realiza como
coeficiente de la acción de Holst. En la carta geométrica declarada,
sea $e^I$ una co-tétrada no degenerada, $\omega^{IJ}$ una conexión,
$\Sigma^{IJ}=e^I\wedge e^J$ y
$F^{IJ}=\dd\omega^{IJ}+\omega^I{}_K\wedge\omega^{KJ}$.
Con $\kappa=8\pi G/c^4$, la acción especializada es
\begin{equation}
\begin{aligned}
S[e,\omega,\Psi]
={}&\frac1{2\kappa c}\int
\Sigma_{IJ}\wedge(\star F^{IJ}+\gamma_{\HMT}^{-1}F^{IJ})\\
&-\frac{\Lambda}{\kappa c}\int\operatorname{vol}_e
+S_m[e,\omega,\Psi].
\end{aligned}
\label{ii:eq:holst}
\end{equation}
La convención de corriente es
$\delta_\omega S_m=-\tfrac12\int\delta\omega^{IJ}\wedge\sigma_{IJ}$.
Para esta variación se fija la conexión en la frontera,
$\delta\omega^{IJ}|_{\partial M}=0$, o se utilizan variaciones de
soporte compacto en el interior. Estas condiciones suficientes anulan
el término de borde que resulta de integrar
$\Sigma_{IJ}\wedge(\star+\gamma_{\HMT}^{-1}I)
D_\omega\delta\omega^{IJ}$.
La variación de \eqref{ii:eq:holst} con esa condición da
\[
D_\omega[(\star+\gamma_{\HMT}^{-1}I)\Sigma^{IJ}]
=\kappa c\,\sigma^{IJ}.
\]
Se obtiene usando $\delta F=D_\omega\delta\omega$, integrando por partes
y haciendo nulo el coeficiente de la variación arbitraria
\cite{ii:holst1996}. La orientación de la frontera y la normalización
de la corriente forman parte de esta fórmula.

En firma lorentziana, $\star^2=-I$ sobre bivectores. La identidad
\[
(\star+\gamma^{-1}I)^{-1}
=\frac{\gamma^2}{1+\gamma^2}(-\star+\gamma^{-1}I)
\]
se comprueba multiplicando ambos miembros: los términos cruzados se
cancelan y queda $(1+\gamma^{-2})I$ antes de normalizar.
Así, el parámetro ya construido interviene en la respuesta torsional,
no únicamente en un valor de tabla.
Las identidades
$D_\omega T^I=F^I{}_J\wedge e^J$ y $D_\omega F^{IJ}=0$
mantienen conectados torsión, curvatura y balance.

La realización geométrica conserva la memoria de las rutas y la
soldadura de frontera. La constante cosmológica \(\Lambda\) de
\eqref{ii:eq:holst} es el coeficiente declarado en la acción; no se
identifica con la representación vectorial \(\Lambda(R)\).
Las identidades siguientes conservan las condiciones de holonomía,
suavidad y no degeneración de la realización.

% BEGIN OWNER /Users/ruben/Documents/ChatGPT/jueces y controles/REVISION_ARGUMENTAL_APERTURAS_CIERRES_20260915/II/manuscript_es/sections/gravedad_memoria_afin.tex
% Artículo VII, recuperación aditiva, revisión 04, 10-09-2026.
% RESULTADO_RECUPERADO: S15, 11c_gravedad_holonomia_rev6.tex:16--126.
% El cociclo canónico c27=(1,18) se conserva como antecedente del cómputo
% de rutas; este fragmento prueba su composición y representación.
% La memoria traslacional finita y la amplitud cosmológica s0 tienen
% tipos distintos. El cociclo no evalúa por sí solo una densidad local.

\subsection{Retorno de fase y representación de la memoria acumulada}
\label{ii:vii:sec:retorno-memoria-gravedad}

La cronología de las rutas enriquecidas distingue el retorno de posición,
el retorno de orientación y el retorno de fase. La publicación observable
identifica ciertos estados de llegada, mientras el registro acumulativo
conserva la secuencia recorrida. La representación geométrica debe
transportar ambos datos conjuntamente.

\subsubsection{Cociclo de inversión y cambio de hoja}
\label{ii:vii:subsec:cociclo-memoria}

Sea \(F_{\rm obs}\) la publicación cronológica del transporte TPK.
Sus bloques canónicos tienen la jerarquía
\[
\begin{array}{c|l}
27&\text{retorno posicional con inversión de orientación},\\
54&\text{retorno posicional y orientacional},\\
108&\text{retorno completo de la fase observable}.
\end{array}
\]
El registro cuenta las inversiones de orientación \(g(a)\) y los
cambios efectivos de hoja \(s(a)\). Para una ruta,
\begin{equation}
 c(\gamma)=\sum_{a\subset\gamma}(g(a),s(a)),\qquad
 c(\gamma_2\gamma_1)=c(\gamma_2)+c(\gamma_1).
 \label{ii:vii:eq:cociclo-ruta}
\end{equation}
El cómputo canónico del bloque de longitud \(27\) publica
\(c_{27}=(1,18)\). La composición conserva esos incrementos en
los cuatro bloques que completan el período observable.
En rutas orientadas hacia adelante el registro es acumulativo;
su extensión al grupoide de rutas utiliza \(\mathbb Z^2\) y
\(c(\bar\gamma)=-c(\gamma)\).

\begin{proposition}[Memoria de la vuelta observable completa]
\label{ii:vii:prop:memoria-vuelta-completa}
Escribiendo \(\mathcal K_{27}=F_{\rm obs}^{27}\), la elevación
al registro, para \(x\in\mathcal Q_{\rm obs}\) y \(z\in\mathbb Z^2\), satisface
\[
 \widetilde{\mathcal K}_{27}(x,z)
   =(\mathcal K_{27}x,z+(1,18)),\qquad
 \widetilde{\mathcal K}_{27}^{\,4}(x,z)
   =(x,z+(4,72)).
\]
Tras \(N\) vueltas completas, el incremento es \(N(4,72)\).
\end{proposition}

\begin{proof}
La proyección observable de \(\mathcal K_{27}\) tiene orden cuatro.
El cómputo canónico anterior fija \(c_{27}=(1,18)\) de manera constante
sobre esa órbita observable. La aditividad de
\eqref{ii:vii:eq:cociclo-ruta} suma los cuatro incrementos y da
\(c_{108}=4(1,18)=(4,72)\).
La repetición aplica la misma ley de concatenación y produce
\(Nc_{108}\). El retorno de la primera coordenada conserva el
incremento de la segunda.
\end{proof}

\subsubsection{Representación afín del cociclo}
\label{ii:vii:subsec:representacion-afin-cociclo}

La realización discreta utiliza un elemento \(R_4\) de orden cuatro en
\(\operatorname{Spin}^+(1,3)\), un vector temporal unitario \(u\)
fijado por su acción vectorial y una unidad de longitud \(\ell_0>0\).
La escala \(\ell_0\) conserva su procedencia en la sección dimensional.
Definimos
\begin{equation}
 \rho_{\rm C}^{\rm disc}(\gamma)=
 \bigl(R_4^{c_g(\gamma)},\,\ell_0c_s(\gamma)u\bigr)
 \in\operatorname{Spin}^+(1,3)\ltimes\mathbb R^{1,3}.
 \label{ii:vii:eq:representacion-discreta-memoria}
\end{equation}
La potencia actúa sobre el vector mediante la representación vectorial
asociada del grupo de espín.

\begin{theorem}[Composición y amplitud traslacional del retorno]
\label{ii:vii:thm:traslacion-memoria}
La aplicación \(\rho_{\rm C}^{\rm disc}\) conserva la identidad,
la concatenación y la inversión. En los retornos completos,
\begin{equation}
 \rho_{\rm C}^{\rm disc}(\gamma_{108})=(I,72\ell_0u),\qquad
 \rho_{\rm C}^{\rm disc}(\gamma_{108}^{\,N})=(I,72N\ell_0u).
 \label{ii:vii:eq:traslacion-retorno-108}
\end{equation}
La inversión cambia el signo de la traslación.
\end{theorem}

\begin{proof}
El producto semidirecto es
\((R,a)(S,b)=(RS,a+Rb)\). Como \(R_4u=u\),
\[
\begin{split}
 \rho_{\rm C}^{\rm disc}(\gamma_2)
 \rho_{\rm C}^{\rm disc}(\gamma_1)
 &=
 \bigl(R_4^{c_g(\gamma_2)+c_g(\gamma_1)},
 \ell_0(c_s(\gamma_2)+c_s(\gamma_1))u\bigr)\\
 &=\rho_{\rm C}^{\rm disc}(\gamma_2\gamma_1).
\end{split}
\]
El cociclo de la identidad es cero y el de la inversa es el opuesto.
Para la vuelta completa se sustituyen \(c_{108}=(4,72)\) y
\(R_4^4=I\). La concatenación de \(N\) vueltas mantiene la parte
rotacional identidad y suma sus traslaciones.
\end{proof}

La amplitud \(72\ell_0\) es una longitud asociada a una ruta finita.
Una realización mediante cotetrada y conexión la transporta a su
holonomía afín. La contribución local a las ecuaciones gravitatorias
se obtiene después mediante la curvatura de esa conexión, la
variación de la acción material y la normalización de su densidad.
La representación conserva así el contenido acumulativo exacto que
debe intervenir en esa realización.

% END OWNER


% BEGIN OWNER /Users/ruben/Documents/ChatGPT/jueces y controles/REVISION_ARGUMENTAL_APERTURAS_CIERRES_20260915/II/manuscript_es/sections/gravedad_pantalla_soldadura.tex
% Artículo VII. Incorporación de resultados preexistentes, 10-09-2026.
% Propietarios íntegros preservados en fuentes_conservadas:
% 17_pantalla_cubica_soldadura.tex (P), 18_clausura_energetica.tex (E).
% P: incidencia, cociente, soldadura celular y refinamiento.
% E: respuesta, unicidad, extensión mínima y expectativa tracial.
% Dependencias anteriores de VII: núcleo APP--TRIT--TPK; realización
% cúbica de la ruta, cociclo de transporte y compresión q_vis=5/9.
% Este fragmento no constituye por sí solo el artículo autónomo.
% La identidad de pantalla se distingue de la contracción material de espín.

\subsection{Incidencia cúbica y soldadura de frontera}
\label{ii:vii:sec:pantalla-respuesta}

La realización cúbica del transporte conserva seis incidencias orientadas.
Su cociente geométrico, la forma lorentziana y la respuesta energética
se obtienen de la misma matriz de incidencias. El transporte procede de
la ruta enriquecida APP--TRIT--TPK: la cara, el signo, el acarreo y la
memoria acompañan a cada arista. En esta sección se desarrolla la
operación lineal que actúa sobre esa publicación de frontera.

\subsubsection{Incidencia cúbica y cociente de rango cuatro}
\label{ii:vii:sec:cociente-pantalla}

Sea \(H_F=\mathbb R^F\), con
\(F=\{(i,\epsilon):1\leq i\leq3,\ \epsilon\in\{-1,1\}\}\),
su producto escalar euclidiano y su base \(e_{i,\epsilon}\).
La oposición de caras es la involución ortogonal
\(Re_{i,\epsilon}=e_{i,-\epsilon}\). Definimos
\[
 u=\sum_{i,\epsilon}e_{i,\epsilon},\qquad
 P_u=\frac{uu^{\mathsf T}}6,\qquad P_-=\frac{I-R}{2},\qquad
 P_0=\frac{I+R}{2}-P_u,\qquad P_\square=P_u+P_-.
\]
Los vértices son \(V=\{-1,1\}^3\). La matriz de incidencia
\(M:\mathbb R^V\longrightarrow H_F\) queda determinada por
\[
 M_{(i,\epsilon),\nu}=\mathbf1_{\{\nu_i=\epsilon\}}.
\]

\begin{proposition}[Descomposición de la incidencia]
\label{ii:vii:prop:incidencia}
Se tiene
\[
 MM^{\mathsf T}=12P_u+4P_-,
 \qquad
 \ker M^{\mathsf T}=P_0H_F.
\]
Por consiguiente, el cociente
\(Q_\square=Q=H_F/P_0H_F\), identificado ortogonalmente con
\(P_\square H_F\), tiene dimensión cuatro.
\end{proposition}

\begin{proof}
Una cara contiene cuatro vértices. Dos caras opuestas carecen de vértices
comunes y dos caras de ejes diferentes tienen dos vértices comunes.
Escribiendo \(\mathbf J=uu^{\mathsf T}\), resulta
\[
 MM^{\mathsf T}=4I+2(\mathbf J-I-R)
 =2(I-R)+2\mathbf J=4P_-+12P_u.
\]
Los tres proyectores \(P_u,P_-,P_0\) son ortogonales y suman \(I\);
sus rangos son, respectivamente, uno, tres y dos. La identidad
\(\|M^{\mathsf T}x\|^2=\langle x,MM^{\mathsf T}x\rangle\)
identifica el núcleo anunciado. El rango complementario es cuatro.
\end{proof}

Sobre \(Q\), la forma
\[
 B=P_--P_u
\]
tiene signatura \((3,1)\). La orientación temporal se fija por \(u\).
Con
\[
 t=\frac{u}{\sqrt6},\qquad
 d_i=\frac{e_{i,+}-e_{i,-}}{\sqrt2},\qquad
 W=\sqrt6P_u+\sqrt2P_-,
\]
se cumple \(We_{i,\epsilon}=t+\epsilon d_i\). Esos seis vectores
son nulos para \(B\), mientras que
\[
 3t+\nu_1d_1+\nu_2d_2+\nu_3d_3
\]
tiene norma cuadrática \(-6\). La misma incidencia satisface
\(MM^{\mathsf T}=2W^*W\) sobre \(Q\). Las rotaciones propias del cubo
conmutan con los proyectores, preservan ambas formas y fijan \(t\).
Estas identidades determinan la métrica de la publicación cúbica.

\subsubsection{Soldadura celular y transporte de memoria}
\label{ii:vii:sec:soldadura-celular}

Para cada arista orientada \(a\), la publicación de la ruta proporciona
la cara \(\lambda(\underline a)\), el signo
\(\sigma(a)\), las fibras de origen y destino y el transporte
\(U(a):Q_{s(a)}\longrightarrow Q_{t(a)}\).
La inversión verifica \(U(\bar a)=U(a)^{-1}\).
Las semiaristas de frontera conservan su inversión formal y su incidencia
en el borde. Con \(\ell_m=\ell_0\,9^{-m}>0\), la soldadura es
\begin{equation}
 \beta_m(a)=\ell_m\sigma_m(a)
              n_{\lambda(\underline a)}^{s(a)},\qquad
 n_{i,\epsilon}=t+\epsilon d_i.
 \label{ii:vii:eq:soldadura-arista}
\end{equation}
La cara y el signo se transportan junto con la memoria de la ruta.
La convención de inversión exige
\begin{equation}
 \beta_m(\bar a)=-U_m(a)\beta_m(a).
 \label{ii:vii:eq:inversion-soldadura}
\end{equation}
Aplicando dos veces esta relación se recupera \(\beta_m(a)\);
el registro de acarreo conserva, por separado, la composición ordenada.

\begin{proposition}[Naturalidad de la soldadura por refinamiento]
\label{ii:vii:prop:refinamiento-soldadura}
Sea \(a=a_1\cdots a_9\) un refinamiento compatible de una arista.
Supóngase que las nueve incidencias hijas, transportadas a la fibra del
padre, conservan la cara y el signo, y que
\(\ell_{m+1}=\ell_m/9\). Sea
\(\mathcal U_j:Q_{m,s(a)}\longrightarrow Q_{m+1,s(a_j)}\)
el transporte acumulado que incluye la identificación por refinamiento
y el trayecto hasta el origen del hijo \(j\). Entonces
\[
 \beta_m(a)=
 \sum_{j=1}^{9}\mathcal U_j^{-1}\beta_{m+1}(a_j).
\]
La igualdad es compatible con la inversión y con la concatenación
del registro de memoria.
\end{proposition}

\begin{proof}
Cada sumando, expresado en la fibra del origen del padre, es
\((\ell_m/9)\sigma_m(a)n_{\lambda(\underline a)}\).
La suma de nueve términos reproduce la soldadura del padre.
La inversión invierte el orden de los transportes y sustituye cada uno
por su inverso; la relación \eqref{ii:vii:eq:inversion-soldadura} proporciona
el signo global. La memoria se concatena en ese mismo orden, por lo
que el refinamiento conserva tanto la publicación vectorial como su
registro enriquecido.
\end{proof}

% END OWNER


% BEGIN OWNER /Users/ruben/Documents/ChatGPT/jueces y controles/REVISION_ARGUMENTAL_APERTURAS_CIERRES_20260915/II/manuscript_es/sections/gravedad_realizacion_geometrica.tex
% RESULTADO_RECUPERADO: S15, 11c_gravedad_holonomia_rev6.tex:128--212;
% S01, 02_dinamica_tres_hojas.tex:155--219; pantalla y respuesta de30.
% Pruebas de curvatura, covariancia y Bianchi reunidas para autonomía.
% La realización suave conserva sus condiciones de holonomía explícitas.

\subsection{Realización geométrica del transporte y de la soldadura}
\label{ii:vii:sec:realizacion-cartan}

La representación discreta del retorno y la soldadura de pantalla
determinan dos aspectos del transporte: la acumulación de memoria en
las rutas y la comparación de sus fibras. Su realización geométrica
conserva ambas operaciones. En esta sección se fijan el espacio de
llegada, la normalización dimensional de la conexión y las identidades
que utilizará su variación.

\subsubsection{Pantalla, cotétrada y conexión}

El cociente de pantalla \(Q_\square\), su forma lorentziana y la
soldadura \(\beta_m\) proceden de
\ref{ii:vii:sec:pantalla-respuesta}. En cada celda, una realización de
pantalla \(J_{{\rm scr},m}\) transporta la forma y la soldadura
a la fibra geométrica correspondiente. Su imagen celular es
\[
 e_m=J_{{\rm scr},m}\beta_m.
\]
La escala \(\ell_m\) permanece dentro de la soldadura. Los
cambios de marco actúan sobre \(J_{{\rm scr},m}\), la conexión
y el transporte de las celdas conjuntamente.

Una realización suave de esos datos sobre una variedad orientada
\(\mathcal U\) de dimensión cuatro consiste en una cotétrada
no degenerada \(e=(e^I)\), una conexión de Lorentz
\(\omega=(\omega^I{}_J)\) y una realización de rutas
\(\mathfrak R_{\rm C}\) que satisfacen
\begin{equation}
 \operatorname{Hol}_{\mathcal A_{\ell_0}}
       (\mathfrak R_{\rm C}(\gamma))
 =\rho_{\rm C}(\operatorname{Hol}_{\rm TPK}(\gamma)).
 \label{ii:vii:eq:compatibilidad-holonomia-geometrica}
\end{equation}
En esta igualdad matricial, la representación se normaliza mediante
\[
 N_{\ell_0}(R,a)=(\Lambda(R),a/\ell_0),\qquad
 \rho_{\rm C}=N_{\ell_0}\circ\rho_{\rm C}^{\rm disc},
\]
donde \(\Lambda:\operatorname{Spin}^+(3,1)\to SO^+(3,1)\)
es la representación vectorial. La aplicación \(N_{\ell_0}\)
preserva el producto semidirecto: dividir la traslación
\(a+\Lambda(R)b\) por \(\ell_0\) equivale a componer las
dos traslaciones normalizadas. Por ello la vuelta completa publica
\(72u\) en la conexión normalizada y \(72\ell_0u\) en la
carta dimensional. El levantamiento de espín conserva, por separado,
\(R_4^{c_g(\gamma)}\); la representación vectorial tiene el
núcleo central \(\{1,-1\}\), de modo que se retiene ese
levantamiento al utilizar campos espinoriales.
Identidad, concatenación, inversión y subdivisión se mantienen en la
realización correspondiente. El emparejamiento de
pantalla y la condición de holonomía especifican los datos de la
realización utilizada en las ecuaciones de campo.

Con \(\eta=\operatorname{diag}(-1,1,1,1)\), se define
\(g=\eta_{IJ}e^I\otimes e^J\). Para una escala elemental
positiva \(\ell_0\), constante en la carta, la conexión afín es
\begin{equation}
 \mathcal A_{\ell_0}
 =\begin{pmatrix}\omega&\ell_0^{-1}e\\0&0\end{pmatrix}.
 \label{ii:vii:eq:conexion-afin-realizada}
\end{equation}
La cotétrada tiene dimensión de longitud; el factor
\(\ell_0^{-1}\) homogeneiza su componente traslacional.

\begin{proposition}[Curvatura y torsión de la conexión realizada]
\label{ii:vii:prop:curvatura-afin}
La curvatura de \eqref{ii:vii:eq:conexion-afin-realizada} es
\begin{equation}
 \mathcal F_{\ell_0}
 =\mathrm d\mathcal A_{\ell_0}
   +\mathcal A_{\ell_0}\wedge\mathcal A_{\ell_0}
 =\begin{pmatrix}F_\omega&\ell_0^{-1}T\\0&0\end{pmatrix},
 \quad F_\omega=\mathrm d\omega+\omega\wedge\omega,
 \quad T=\mathrm de+\omega\wedge e.
 \label{ii:vii:eq:curvatura-afin-realizada}
\end{equation}
\end{proposition}
\begin{proof}
El producto de matrices de formas produce en el bloque diagonal
\(\omega\wedge\omega\) y en el traslacional
\(\ell_0^{-1}\omega\wedge e\). La derivada exterior aporta
\(\mathrm d\omega\) y \(\ell_0^{-1}\mathrm de\), pues
\(\ell_0\) es constante. La suma demuestra la igualdad.
\end{proof}

La curvatura describe el defecto de transporte de marcos y la torsión
el de la cotétrada. Ambas pertenecen a la misma conexión realizada;
sus componentes conservan dominios tensoriales distintos.

\begin{proposition}[Covariancia e identidades de Bianchi]
\label{ii:vii:prop:bianchi-cartan}
Bajo un cambio de marco \(g_L:\mathcal U\to SO^+(3,1)\),
\[
 e'=g_Le,\qquad
 \omega'=g_L\omega g_L^{-1}-\mathrm dg_L\,g_L^{-1},
\]
se cumple \(T'=g_LT\) y \(F_{\omega'}=g_LF_\omega g_L^{-1}\).
Además,
\begin{equation}
 D_\omega T=F_\omega\wedge e,\qquad D_\omega F_\omega=0.
 \label{ii:vii:eq:bianchi-realizacion}
\end{equation}
\end{proposition}
\begin{proof}
En \(\mathrm d(g_Le)+\omega'\wedge g_Le\), los términos
que contienen \(\mathrm dg_L\wedge e\) se cancelan.
La misma sustitución en la curvatura da su transformación por
conjugación. Para la primera identidad de Bianchi, se expande
\[
 D_\omega(\mathrm de+\omega\wedge e)
 =(\mathrm d\omega+\omega\wedge\omega)\wedge e.
\]
En \(\mathrm dF_\omega+\omega\wedge F_\omega-F_\omega\wedge\omega\),
los términos que contienen \(\mathrm d\omega\) y los productos
cúbicos se cancelan por pares. Esto prueba la segunda identidad.
\end{proof}

\subsubsection{Las tres realizaciones cuadráticas de curvatura}

La hoja \(\tau\in\{+1,0,-1\}\) del TRIT se representa en
el álgebra de Cartan mediante generadores \(J_{ab}=-J_{ba}\)
y \(P_a^{(\tau)}\), con relaciones
\begin{align*}
 [J_{ab},J_{cd}]&=
 \eta_{bc}J_{ad}-\eta_{ac}J_{bd}
 -\eta_{bd}J_{ac}+\eta_{ad}J_{bc},\\
 [J_{ab},P_c^{(\tau)}]&=
 \eta_{bc}P_a^{(\tau)}-\eta_{ac}P_b^{(\tau)},\\
 [P_a^{(\tau)},P_b^{(\tau)}]&=-\tau J_{ab}.
\end{align*}
Para \(\ell>0\) constante, se escribe
\[
 \mathcal A_\tau^{\rm Cart}
 =\tfrac12\omega^{ab}J_{ab}+\ell^{-1}e^aP_a^{(\tau)}.
\]

\begin{proposition}[Curvatura de la realización tritificada]
\label{ii:vii:prop:curvatura-tritificada}
Con las relaciones anteriores,
\begin{equation}
 \mathcal F_\tau^{\rm Cart}
 =\tfrac12(F_\omega^{ab}-\tau\ell^{-2}e^a\wedge e^b)J_{ab}
     +\ell^{-1}T^aP_a^{(\tau)}.
 \label{ii:vii:eq:curvatura-tritificada}
\end{equation}
\end{proposition}
\begin{proof}
Se expande \(\mathrm d\mathcal A+\tfrac12[\mathcal A,\mathcal A]\).
El corchete del sector de Lorentz produce \(F_\omega\), el
mixto produce \(D_\omega e\), y el traslacional aporta
\(-\tfrac12\tau\ell^{-2}e^a\wedge e^bJ_{ab}\).
\end{proof}

En la hoja central se recupera la conexión afín; las hojas extremas
conservan los dos signos del término de curvatura constante. La
representación mantiene así el régimen trítico de la ruta antes de
aplicar las ecuaciones variacionales. La igualdad de holonomía
\eqref{ii:vii:eq:compatibilidad-holonomia-geometrica} corresponde a la
hoja central. En cada hoja extrema se utiliza la representación
\(\rho_{{\rm C},\tau}\) en su grupo de Cartan y su condición
de holonomía respectiva. La corriente y el tensor métrico utilizan
los campos y la representación de la misma realización.

% END OWNER


\subsection{Transporte orientado de la conexión y del área}
\label{ii:sec:ii-costura-barbero}
La paridad del funcional de incidencia permite comparar las cartas
\(\gamma\) y \(-\gamma\). En la realización canónica geométrica,
sea \(E^a_i\) una triada densitizada, \(\Gamma_a^i(E)\) su
conexión intrínseca y \(K_a^i\) la curvatura extrínseca en la
convención normal elegida. La conexión de Ashtekar--Barbero es
\(\mathcal A_a^i=\Gamma_a^i(E)+\gamma K_a^i\).

\begin{proposition}[Involución entre cartas orientadas]
\label{ii:prop:ii-costura-barbero}
Para \(\gamma\in\mathbb R\setminus\{0\}\), la aplicación
\[
 \mathfrak b:(\gamma,K,E)\longmapsto(-\gamma,-K,E)
\]
es involutiva. Conserva \(\mathcal A\) y revierte la forma
\[
 \Omega_{KE}=\frac1\kappa\int_\Sigma
 \delta K_a^i\wedge\delta E^a_i,
\]
con \(\kappa\) fijo y las mismas condiciones de frontera.
\end{proposition}
\begin{proof}
Aplicar dos veces \(\mathfrak b\) restituye las tres coordenadas.
Como \(E\) permanece fijo, también permanece fija \(\Gamma(E)\).
El producto de las dos coordenadas impares satisface
\((-\gamma)(-K)=\gamma K\), y conserva la conexión.
En el espacio de variaciones, \(\delta(-K)=-\delta K\) y
\(\delta E\) permanece fijo: el pullback de la forma es
\(\mathfrak b^*\Omega_{KE}=-\Omega_{KE}\).
\end{proof}

El área orientada \(\mathcal A_{\rm or}=\gamma a_0N\) y su
magnitud física \(\mathcal A_{\rm fis}=|\gamma|a_0N\), para
\(a_0>0\) y \(N\geq0\), coinciden en la carta positiva.
Al prolongar a la carta opuesta, la primera es impar y la segunda
permanece invariante. Esta distinción conserva tanto la orientación
de la incidencia como la positividad del observable geométrico.
La proposición establece el transporte algebraico de estas cartas;
su realización como inversión de la normal espacial utiliza,
además, el transporte de la materia y de los datos de frontera
de la acción~\eqref{ii:eq:holst}.
